\documentclass[10pt]{beamer}

%%%
% PREAMBLE FOR THIS DOC 
%%%
%https://tex.stackexchange.com/questions/68821/is-it-possible-to-create-a-latex-preamble-header
\usepackage{/Users/miw267/Repos/latex_preamble/beamer_preamble}

\usetikzlibrary{matrix}



% \album{image file}{url}{artist}{album}
\newcommand{\album}[4]{%
  \begin{minipage}[t]{0.185\textwidth}\centering
    \href{#2}{\includegraphics[width=\linewidth]{#1}}\\[0.3em]
    {\small\bfseries #3}\\{\small #4}
  \end{minipage}}
  
\newcommand{\albumtiny}[4]{%
  \begin{minipage}[t]{0.185\textwidth}\centering
    \href{#2}{\includegraphics[width=\linewidth]{#1}}\\[0.3em]
    {\tiny \bfseries #3}\\{\tiny #4}
  \end{minipage}}
  
  
%%%
% DOCUMENT
%%%

\begin{document}

%\maketitle

%% Title page frame
%\begin{frame}
%    \titlepage 
%\end{frame}





\title{10/05/2026: Midterm Review Part 2: More Proofs}
\author{CSCI 246: Discrete Structures}

\begin{frame}
    \titlepage 
\end{frame}


\begin{frame}

\begin{myredbox}[title=Announcements]
\begin{itemize}
	\item Let's review some midterm logistics.
	\item If you prepared for today's optional quiz retake, you're 1/5 of the way done studying for the midterm!  Yippee!
\end{itemize}
\end{myredbox}
\vfill 


\begin{myyellowbox}[title=Today's Agenda]
\begin{itemize}	
	
	\item Midterm discussion (5 mins)
	\item Optional Quiz Retake (20 mins)
	\item Group exercises (25 mins)
%	%
%	\begin{itemize}
%	\footnotesize 
%	\item Review induction 
%	\end{itemize}
%	%
\end{itemize}

%	Rationale for group exercises: we got shortchanged on time last couple days, and I already did a lot of lectures, so I want you to practice. Next problems quiz will cover relations and functions: Hamkins and 
%	
\end{myyellowbox}
\vfill 

\end{frame}



\begin{frame}[standout]
Group exercises
\end{frame}

\begin{frame}
\begin{columns}
\begin{column}{0.33\textwidth}
Ashley, Isabelle: 7 \\ 
Bailey, Tanna: 7 \\ 
Blaisdell, Olyver: 5 \\ 
Brown, Ashton: 1 \\ 
Brown, Jonathan: 2 \\ 
Cardenas, Dominic: 2 \\ 
Cutler, George: 11 \\ 
Fabik, Roland: 8 \\ 
Fellows, Audrey: 10 \\ 
Fitzgerald, Ryan: 8 \\ 
Fuller, Ryan: 6 \\ 
Griffen, Clara: 4 \\ 
Grutsch, Tanner: 12 \\\end{column}
\begin{column}{0.33\textwidth}
Hager, Nathan: 11 \\ 
Horner, Jayden: 7 \\ 
Johnson, Jaret: 3 \\ 
Kintzel, Samantha: 1 \\ 
Lameres, Alexis: 9 \\ 
Lee, Holden: 5 \\ 
Mayala, Reece: 6 \\ 
McLean, Connor: 12 \\ 
Mcbride, Daniel: 10 \\ 
Miller, Maverick: 2 \\ 
Moss, Emily: 8 \\ 
Mousseau, Matthew: 3 \\ 
Neuman, Evan: 13 \\\end{column}
\begin{column}{0.33\textwidth}
Putnam, John: 5 \\ 
Ramos, Vydel: 10 \\ 
Rathbun, Jedidiah: 9 \\ 
Roberts, Matthew: 13 \\ 
Roller, Nathan: 4 \\ 
Ryan, Otis: 3 \\ 
Shevtsov, Timothy: 1 \\ 
Sinclair, Aidan: 11 \\ 
Smith, Charles: 6 \\ 
Stensland, Emma: 13 \\ 
Thiessen, Logan: 9 \\ 
Thompson, Jaiden: 4 \\ 
Vath, Ethan: 12 \\\end{column}
\end{columns}
\end{frame}


\begin{frame}{Group exercises}
\footnotesize
\textbf{New group exercises on proofs.}
\vspace{-.3em}
\begin{enumerate}
	\item  Prove that the product of two even integers is even.
	\item  Let $x$ be an integer. Prove that $0 \mid x$ if and only if $x=0$.
	\item Disprove: If $p$ and $q$ are prime, then $p+q$ is composite.
	\item Let $A = \set{x \in \mathbb{Z} : 4 \mid x}$ and let $B = \set{x \in \mathbb{Z} : 2 \mid x}$. Show that $A \subseteq B$. 
	\item Let $A = \set{n \in \mathbb{Z} : 3 \mid n}$ and let $B = \set{n \in \mathbb{Z} : 3 \mid (n+6)}$. Show that $A=B$. 
\end{enumerate}
\vfill 
\vspace{-.5em}
\textbf{Practicing some of the older group exercises on proofs.}
\vspace{-.3em}
\begin{enumerate}
\item A \textbf{tautology} is a Boolean expression that evaluates to \texttt{TRUE} for all possible values of its variables. Use truth tables to show the following is a tautology:
\[ (x \implies y) \land (y \implies z) \implies (x \implies z) \]
%\item For each of the following relations defined on the integers, please state whether the relation is reflexive, symmetric, and/or transitive:  (a) $\leq$ (less than or equal to), (b) $\mid$ (divides). Please justify your answer for part (b).
\item Let $n$ be a positive integer. Prove that congruence modulo $n$ is an equivalence relation on the set of integers.
\item The distributive properties for sets are
\[ A \cup (B \cap C) = (A \cup B) \cap (A \cup C), \qquad \text{and} \qquad A \cap (B \cup C) = (A \cap B) \cup (A \cap C). \]
Prove the second property.
\end{enumerate}
%\vfill
%\small 
%\begin{mygreenbox}[title={Reference: Definitions}]
%\begin{itemize}
%	\item Let $a,b$ be integers.  We say that $a$ is \textbf{divisible} by $b$, written $b \mid a$, if there is an integer $c$ such that $bc=a$. We call $b$ a \textbf{divisor} of $a$.
%	\item An integer is called \textbf{even} if it is divisible by 2. (That is, an integer $x$ is even if there is an integer $c$ such that $x=2c$.)
%	\item An integer $p$ is called \textbf{prime} if $p>1$ and the only positive divisors of $p$ are 1 and $p$.
%	\item A positive integer $a$ is called \textbf{composite} if there is an integer $b$ such that  $1<b<a$ and $b \mid a$.
%\end{itemize}
%	
%\end{mygreenbox}

\end{frame}





\begin{frame}{Solution: Group exercise \#1}
\textbf{Problem.} Prove that the product of two even integers is even.
\vfill 
\textbf{Solution.}   We show that if $z=xy$ where $x$ and $y$ are even, then $z$ is even.   Let $x$ and $y$ be even.  Then $x=2a$ where $a$ is an integer and $y=2b$ where $b$ is an integer.  So $z=xy=(2a)(2b)=(4ab)=2(2ab)$. Hence $z=2c$ for some integer $c$ (namely $c= 2ab$).  Hence $z$ is even. 

\vfill
\begin{mygreenbox}[title={Reference: Definition}]
An integer is called \textbf{even} if it is divisible by 2. (That is, an integer $z$ is even if there is an integer $c$ such that $z=2c$.)
\end{mygreenbox}

\end{frame}

\begin{frame}{Solution: Group exercise \#2}
\textbf{Problem.} Let $x$ be an integer. Prove that $0 \mid x$ if and only if $x=0$.
\vfill 
\textbf{Solution.}  First we show that if $0 \mid x$ then $x=0$.  Let $0 \mid x$. Then there is an integer $c$ such that $0c=x$.  So $x=0$.

Now we show that if $x=0$ then $0 \mid x$.  Let $x=0$.  Let $c$ be any integer (e.g. $c=7$).  Then $x=0c$. Hence, there is an integer $c$ such that $0c=x$.  Hence $0 \mid x$. 

\vfill
\begin{mygreenbox}[title={Reference: Definition}]
Let $a,b$ be integers.  We say that $a$ is \textbf{divisible} by $b$, written $b \mid a$, if there is an integer $c$ such that $bc=a$. 
\end{mygreenbox}

\end{frame}



\begin{frame}{Solution: Group exercise \#3 }
\textbf{Problem.} Disprove: If $p$ and $q$ are prime, then $p+q$ is composite.
	\vfill
\textbf{Solution.} We give a counterexample where $p$ and $q$ are prime, but $p+q$ is not composite.  Let $p=2$ and $q=3$.  Then $p$ and $q$ are both prime. ($2>1$ and the only positive divisors of $2$ are 1 and $2$, and  $3>1$ and the only positive divisors of $3$ are 1 and $3$.). However, $p+q=5$, which is prime  ($5>1$ and the only positive divisors of $5$ are 1 and $5$), and therefore not composite.
\vfill 
\begin{mygreenbox}[title={Reference: Definitions}]
\begin{itemize}
	\item An integer $p$ is called \textbf{prime} if $p>1$ and the only positive divisors of $p$ are 1 and $p$.
	\item A positive integer $a$ is called \textbf{composite} if there is an integer $b$ such that  $1<b<a$ and $b \mid a$.
\end{itemize}
	
\end{mygreenbox}

\end{frame}


\begin{frame}{Solution: Group exercise \#4}
\textbf{Problem.}   Let $A = \set{x \in \mathbb{Z} : 4 \mid x}$ and let $B = \set{x \in \mathbb{Z} : 2 \mid x}$. Show that $A \subseteq B$. 
\vfill
\textbf{Solution.} We will show that if $x \in A$, then $x \in B$. Let $x \in A$.  Then $4 \mid x$. So there is an integer $c$ such that $x=4c$.  So $x=2(2c)$.  So $x=2d$ for some integer $d$ (namely $d=2c$). So $2 \mid x$.  So $x \in B$.

\vfill
\begin{mygreenbox}[title={Reference: Definition}]
Let $a,b$ be integers.  We say that $a$ is \textbf{divisible} by $b$, written $b \mid a$, if there is an integer $c$ such that $bc=a$. We call $b$ a \textbf{divisor} of $a$.	
\end{mygreenbox}

\end{frame}




\begin{frame}{Solution: Group exercise \#5}
\textbf{Problem.}   Let $A = \set{n \in \mathbb{Z} : 3 \mid n}$ and let $B = \set{n \in \mathbb{Z} : 3 \mid (n+6)}$. Show that $A=B$. 
\vfill
\textbf{Solution.} First we show that if $x \in A$, then $x \in B$. Let $x \in A$.  Then $3 \mid x$.   So $x=3c$ for some integer $c$. So $x+6=3c+6$. So $x+6=3(c+2)$. So $x+6=3d$ for some integer $d$, namely $d=c+2$.   So $3 \mid (x+6)$. So $x \in B$.

Now we show that if $x \in B$, then $x \in A$. Let $x \in B$. So $3 \mid (x+6)$.  So $3c=x+6$ for some integer $c$.   So $3c-6=x$. So $3(c-2)=x$.  So $3d=x$ for some integer $d$ (namely $d=c-2$).   So $x \in A$. 
\vfill
\begin{mygreenbox}[title={Reference: Definition}]
Let $a,b$ be integers.  We say that $a$ is \textbf{divisible} by $b$, written $b \mid a$, if there is an integer $c$ such that $bc=a$. We call $b$ a \textbf{divisor} of $a$.	
\end{mygreenbox}

\end{frame}


\end{document}
